\documentclass[11pt]{article}

\usepackage[a4paper,margin=2.5cm]{geometry}
\usepackage{amsmath,amssymb,amsthm,amsfonts}
\usepackage{mathtools}
\usepackage{hyperref}
\usepackage{graphicx}

\title{Measure Preservation and Dyadic Volume Decomposition\\
	in the EABC Coordinate Space}

\author{Thomas Hoffbauer}

\date{\today}

\newtheorem{theorem}{Theorem}
\newtheorem{lemma}{Lemma}
\newtheorem{definition}{Definition}
\newtheorem{corollary}{Corollary}
\newtheorem{remark}{Remark}

\begin{document}
	
	\maketitle
	
	\begin{abstract}
		
		We investigate an infinite dyadic decomposition of a local source volume
		within the four-dimensional EABC coordinate system.
		The construction induces a natural probability measure on the four
		orthogonal coordinate directions
		
		\[
		(E,A,B,C).
		\]
		
		We prove that the decomposition is measure preserving,
		convergent, and complete.
		The resulting weight vector
		
		\[
		\left(\frac{8}{15},\frac{4}{15},\frac{2}{15},\frac{1}{15}\right)
		\]
		
		defines a canonical point in the positive simplex and simultaneously
		determines a normalized quaternionic direction on the unit
		three-sphere.
		
		The decomposition therefore separates additive volume conservation
		from geometric direction normalization.
		
	\end{abstract}
	
	\section{Introduction}
	
	The EABC framework is based on the four invertible residue classes
	modulo $12$
	
	\[
	E\equiv1,\qquad
	A\equiv5,\qquad
	B\equiv7,\qquad
	C\equiv11
	\pmod{12}.
	\]
	
	These classes form the Klein four-group
	
	\[
	V_4=\{E,A,B,C\}
	\]
	
	with multiplication laws
	
	\[
	A^2=B^2=C^2=E,
	\qquad
	AB=C,
	\qquad
	AC=B,
	\qquad
	BC=A.
	\]
	
	In the present work we consider a local source volume
	spanned by four orthogonal coordinate directions
	
	\[
	\hat e,\hat a,\hat b,\hat c.
	\]
	
	A dyadic four-phase decomposition generates a canonical hierarchy
	of weights associated with these directions.
	
	\section{Dyadic Four-Phase Expansion}
	
	\begin{definition}
		
		Define the EABC expansion
		
		\[
		\Xi
		=
		\sum_{k=0}^{\infty}
		\left(
		2^{-(4k+1)}\hat e
		+
		2^{-(4k+2)}\hat a
		+
		2^{-(4k+3)}\hat b
		+
		2^{-(4k+4)}\hat c
		\right).
		\]
		
	\end{definition}
	
	Since the EABC cycle repeats every four steps,
	each coordinate direction is scaled by
	
	\[
	q=2^{-4}=\frac1{16}.
	\]
	
	The individual coordinate contributions are therefore geometric series.
	
	\section{Coordinate Weights}
	
	\begin{lemma}
		
		The four coordinate weights are
		
		\[
		\Xi_e=\frac{8}{15},
		\qquad
		\Xi_a=\frac{4}{15},
		\qquad
		\Xi_b=\frac{2}{15},
		\qquad
		\Xi_c=\frac{1}{15}.
		\]
		
	\end{lemma}
	
	\begin{proof}
		
		For the $E$ direction,
		
		\[
		\Xi_e
		=
		\sum_{k=0}^{\infty}
		2^{-(4k+1)}
		=
		\frac12
		\sum_{k=0}^{\infty}
		\left(\frac1{16}\right)^k.
		\]
		
		Using the geometric-series identity,
		
		\[
		\sum_{k=0}^{\infty}q^k
		=
		\frac1{1-q},
		\qquad |q|<1,
		\]
		
		gives
		
		\[
		\Xi_e
		=
		\frac12
		\cdot
		\frac1{1-\frac1{16}}
		=
		\frac8{15}.
		\]
		
		The remaining coordinates follow analogously:
		
		\[
		\Xi_a=\frac4{15},
		\qquad
		\Xi_b=\frac2{15},
		\qquad
		\Xi_c=\frac1{15}.
		\]
		
	\end{proof}
	
	\section{Measure Preservation}
	
	\begin{theorem}[EABC Measure Preservation]
		
		The dyadic four-phase decomposition induces a probability measure
		
		\[
		\mu=(\Xi_e,\Xi_a,\Xi_b,\Xi_c)
		\]
		
		satisfying
		
		\[
		\mu(E)+\mu(A)+\mu(B)+\mu(C)=1.
		\]
		
	\end{theorem}
	
	\begin{proof}
		
		Using the previous lemma,
		
		\[
		\frac8{15}
		+
		\frac4{15}
		+
		\frac2{15}
		+
		\frac1{15}
		=
		\frac{15}{15}
		=
		1.
		\]
		
	\end{proof}
	
	Consequently,
	
	\[
	\mu
	=
	\left(
	\frac8{15},
	\frac4{15},
	\frac2{15},
	\frac1{15}
	\right)
	\]
	
	defines a probability measure on the EABC coordinate set.
	
	\section{Volume Conservation}
	
	\begin{definition}
		
		Let the local source volume be normalized by
		
		\[
		V_M=1.
		\]
		
	\end{definition}
	
	\begin{corollary}
		
		The dyadic decomposition preserves total volume in the additive sense:
		
		\[
		V_M
		=
		\Xi_e+\Xi_a+\Xi_b+\Xi_c.
		\]
		
	\end{corollary}
	
	Thus the decomposition neither creates nor destroys volume.
	No gaps and no overlaps occur with respect to the induced measure.
	
	\section{Euclidean Geometry}
	
	The corresponding EABC vector is
	
	\[
	\Xi
	=
	\frac8{15}\hat e
	+
	\frac4{15}\hat a
	+
	\frac2{15}\hat b
	+
	\frac1{15}\hat c.
	\]
	
	Because the coordinate axes are orthogonal,
	
	\[
	\|\Xi\|_2^2
	=
	\left(\frac8{15}\right)^2
	+
	\left(\frac4{15}\right)^2
	+
	\left(\frac2{15}\right)^2
	+
	\left(\frac1{15}\right)^2.
	\]
	
	Hence
	
	\[
	\|\Xi\|_2^2
	=
	\frac{64+16+4+1}{225}
	=
	\frac{85}{225}
	=
	\frac{17}{45}.
	\]
	
	\section{Quaternionic Interpretation}
	
	Identify the EABC basis with the quaternion basis
	
	\[
	E\leftrightarrow1,
	\qquad
	A\leftrightarrow i,
	\qquad
	B\leftrightarrow j,
	\qquad
	C\leftrightarrow k.
	\]
	
	Then
	
	\[
	Q
	=
	\frac8{15}
	+
	\frac4{15}i
	+
	\frac2{15}j
	+
	\frac1{15}k.
	\]
	
	The quaternion norm satisfies
	
	\[
	|Q|^2
	=
	\frac{17}{45}.
	\]
	
	Therefore
	
	\[
	\widehat Q
	=
	\sqrt{\frac{45}{17}}\,Q
	\]
	
	lies on the quaternionic unit sphere
	
	\[
	S^3.
	\]
	
	\section{The Role of Fifteen}
	
	The number
	
	\[
	15=8+4+2+1
	\]
	
	appears naturally as the denominator of the normalized measure.
	
	Remarkably,
	
	\[
	15=2^4-1
	\]
	
	is also the number of non-empty subsets of a four-element set.
	
	Hence the denominator coincides with the complete combinatorial content
	of a four-phase system.
	
	\section{Conclusion}
	
	The EABC dyadic decomposition induces a canonical probability measure
	
	\[
	\left(
	\frac8{15},
	\frac4{15},
	\frac2{15},
	\frac1{15}
	\right)
	\]
	
	that preserves total volume exactly.
	
	The construction separates two distinct structures:
	
	\begin{enumerate}
		\item additive measure conservation,
		\item geometric normalization on the quaternionic unit sphere.
	\end{enumerate}
	
	This provides a mathematically consistent bridge between
	dyadic hierarchical decompositions,
	measure-preserving dynamics,
	and quaternionic geometry.
	
\end{document}